% ── 6. Cloud-Native Application Design (Twelve-Factor App) ───────────────────
\section{Cloud-Native Application Design}

\subsection{Codebase}

The cloud-side platform aligns with the codebase factor by treating the backend services and infrastructure definitions as version-controlled artifacts that support multiple deployments. This is important because the same system must be promoted across development, staging, and production without creating drift between environments. The tradeoff is that this approach depends on disciplined source control and release practices, but it makes changes easier to track, review, and reproduce.

\subsection{Config}

Configuration is kept separate from application logic so that environment-specific values can change without requiring code changes. This is especially important for a platform that will operate across multiple environments and may evolve over time as endpoints, secrets, and service identifiers change. The tradeoff is that externalized configuration improves flexibility and security, but it also increases the need for careful configuration management and validation.

\subsection{Backing Services}

The architecture fits the backing services principle by treating databases, storage, messaging, and workflow services as attached resources rather than tightly embedding them into the application layer. This is a strong fit for the kata because the platform depends on managed services for telemetry, user data, automation, and archived content. The tradeoff is that this improves modularity and maintainability, but it also increases dependence on AWS-managed service patterns.

\subsection{Build, Release, Run}

The design also aligns well with build, release, and run as separate concerns. This matters because a consumer platform controlling household devices should support repeatable deployments and safe rollback when changes introduce faults. The tradeoff is that a cleaner release process improves reliability, but it also requires stronger CI/CD discipline and more structured deployment workflows.

\subsection{Processes and Concurrency}

The backend is designed around stateless compute and horizontal scaling, which makes the processes and concurrency factors highly relevant. This supports the expected growth to thousands of households by allowing the platform to handle more requests and device events through scale-out behavior instead of relying on larger individual servers. The tradeoff is that stateless scaling improves elasticity, but it pushes more responsibility onto external services for state management and coordination.

\subsection{Dev/Prod Parity}

Dev/prod parity is important because differences between environments can easily hide defects until they reach production. The architecture supports this by keeping environments structurally similar and by defining infrastructure in code rather than rebuilding it manually for each stage. The tradeoff is that closer parity reduces deployment surprises, but it can increase development cost compared to using lighter local substitutes.

\subsection{Logs and Admin Processes}

Logs and admin processes are also strong fits for this architecture. Centralized logs are important because the platform spans multiple managed services and asynchronous workflows, while one-off administrative tasks such as cleanup, maintenance, and reporting should remain separate from the main request path. The tradeoff is that this separation improves observability and operational control, but it adds some overhead in monitoring, automation, and access management.

Some Twelve-Factor ideas are less central here. Port binding is only partially applicable because managed entry points such as API Gateway expose the application on behalf of the backend, and disposability is largely handled by the serverless platform rather than by long-running application processes. For that reason, the most relevant factors for this design are the ones tied to deployment consistency, externalized state, attached services, and operational separation.